\documentclass[10pt,a4paper]{amsart}
\usepackage{amsmath,amssymb,amsthm,graphicx}
% Retained minimal preamble for the Stacks Project chapter more-morphisms.tex (More on Morphisms).
% Verbatim excerpt of the chapter preamble: Xy-pic setup (native preamble lines 14-18) and the native
% theorem declarations with the mathematical macros (native preamble lines 154-217). Copied mechanically;
% nothing added, removed or redefined. The wrapper class and the amsmath/amssymb/amsthm/graphicx base
% packages are supplied by the assembler, so the class/externaldocument part of the native preamble is omitted.
% For commutative diagrams we use Xy-pic
\usepackage[all]{xy}

% We use 2cell for 2-commutative diagrams.
\xyoption{2cell}
% Theorem environments.
%
\theoremstyle{plain}
\newtheorem{theorem}[subsection]{Theorem}
\newtheorem{proposition}[subsection]{Proposition}
\newtheorem{lemma}[subsection]{Lemma}

\theoremstyle{definition}
\newtheorem{definition}[subsection]{Definition}
\newtheorem{example}[subsection]{Example}
\newtheorem{exercise}[subsection]{Exercise}
\newtheorem{situation}[subsection]{Situation}

\theoremstyle{remark}
\newtheorem{remark}[subsection]{Remark}
\newtheorem{remarks}[subsection]{Remarks}

\numberwithin{equation}{subsection}

% Macros
%
\def\lim{\mathop{\mathrm{lim}}\nolimits}
\def\colim{\mathop{\mathrm{colim}}\nolimits}
\def\Spec{\mathop{\mathrm{Spec}}}
\def\Hom{\mathop{\mathrm{Hom}}\nolimits}
\def\Ext{\mathop{\mathrm{Ext}}\nolimits}
\def\SheafHom{\mathop{\mathcal{H}\!\mathit{om}}\nolimits}
\def\SheafExt{\mathop{\mathcal{E}\!\mathit{xt}}\nolimits}
\def\Sch{\mathit{Sch}}
\def\Mor{\mathop{\mathrm{Mor}}\nolimits}
\def\Ob{\mathop{\mathrm{Ob}}\nolimits}
\def\Sh{\mathop{\mathit{Sh}}\nolimits}
\def\NL{\mathop{N\!L}\nolimits}
\def\CH{\mathop{\mathrm{CH}}\nolimits}
\def\proetale{{pro\text{-}\acute{e}tale}}
\def\etale{{\acute{e}tale}}
\def\QCoh{\mathit{QCoh}}
\def\Ker{\mathop{\mathrm{Ker}}}
\def\Im{\mathop{\mathrm{Im}}}
\def\Coker{\mathop{\mathrm{Coker}}}
\def\Coim{\mathop{\mathrm{Coim}}}

% Boxtimes
%
\DeclareMathSymbol{\boxtimes}{\mathbin}{AMSa}{"02}

%
% Macros for moduli stacks/spaces
%
\def\QCohstack{\mathcal{QC}\!\mathit{oh}}
\def\Cohstack{\mathcal{C}\!\mathit{oh}}
\def\Spacesstack{\mathcal{S}\!\mathit{paces}}
\def\Quotfunctor{\mathrm{Quot}}
\def\Hilbfunctor{\mathrm{Hilb}}
\def\Curvesstack{\mathcal{C}\!\mathit{urves}}
\def\Polarizedstack{\mathcal{P}\!\mathit{olarized}}
\def\Complexesstack{\mathcal{C}\!\mathit{omplexes}}
% \Pic is the operator that assigns to X its picard group, usage \Pic(X)
% \Picardstack_{X/B} denotes the Picard stack of X over B
% \Picardfunctor_{X/B} denotes the Picard functor of X over B
\def\Pic{\mathop{\mathrm{Pic}}\nolimits}
\def\Picardstack{\mathcal{P}\!\mathit{ic}}
\def\Picardfunctor{\mathrm{Pic}}
\def\Deformationcategory{\mathcal{D}\!\mathit{ef}}
\usepackage{cleveref}
\crefname{originallemma}{original lemma}{original lemmas}
\Crefname{originallemma}{Original lemma}{Original lemmas}
\makeatletter
\newlabel{coherent-lemma-quasi-coherent-affine-cohomology-zero}{{01XB}{source}{}{sourcepointer-1}{}}
\newlabel{coherent-lemma-quasi-coherent-affine-cohomology-zero@cref}{{[originallemma][0][]01XB}{[1][0][]source}{}{}{}}
\makeatother
\crefname{originallemma}{original lemma}{original lemmas}
\Crefname{originallemma}{Original lemma}{Original lemmas}
\makeatletter
\newlabel{schemes-lemma-morphism-into-affine}{{01I1}{source}{}{sourcepointer-2}{}}
\newlabel{schemes-lemma-morphism-into-affine@cref}{{[originallemma][0][]01I1}{[1][0][]source}{}{}{}}
\makeatother
\title{First order thickenings are exactly square-zero quasi-coherent ideal extensions of the structure sheaf}
\author{The Stacks Project Authors}
\date{stacks-project:more-morphisms.tex:local-snapshot-20261011; GNU FDL 1.2 or later}
\begin{document}
\maketitle
\noindent\textit{Editorial note.} All delivered mathematical text is the author's own, copied from the pinned Stacks Project chapter file more-morphisms.tex; no statement or proof step was rewritten, reordered or removed.\par\medskip
\noindent\textit{Editorial note.} The preamble is the chapter's own retained minimal preamble (sources/stacks-more-morphisms/retained-preamble.tex), a verbatim excerpt of the native Xy-pic setup and of the native theorem declarations and mathematical macros; the Xy-pic xy environment is required because the counted proof displays two commutative diagrams. No macro was invented or redefined.\par\medskip
\noindent\textit{Editorial note.} The two cross-chapter standard prerequisites are located in the project's own tag list staged in the snapshot (sources/stacks-chow/tags.txt), at the verbatim lines carrying their canonical tags.\par\medskip
\noindent\textit{Allowed standard prerequisites (editorial).}\begin{itemize}
\item Thickening and first order thickening of schemes over a base. Same chapter, Definition environment labelled definition-thickening (source line 36); the statement refers to 'a first order thickening over S' and to \$X' = (X, A)\$. Supplies the definition that the counted proof has to verify for the locally ringed space (X, A); the verification itself stays inside the counted span.
\item Cohomology of a quasi-coherent sheaf on an affine scheme vanishes in positive degrees. Cohomology of Schemes, Lemma labelled coherent-lemma-quasi-coherent-affine-cohomology-zero (canonical Stacks tag 01XB); invoked at proof lines 114-115. Supplies H\textasciicircum{}1(U, I) = 0, which is what makes A -> B = Gamma(U, O\_X) surjective on an affine open; the rest of the argument is retained in the counted span.
\item Morphisms from a locally ringed space to an affine scheme. Schemes, Lemma labelled schemes-lemma-morphism-into-affine (canonical Stacks tag 01I1); invoked at proof lines 118-120. Supplies the canonical morphism (U, A|\_U) -> Spec(A) coming from the ring map B -> Gamma(U, A); the comparison of local rings is retained in the counted span.
\end{itemize}\medskip
\ifdefined\hypertarget\hypertarget{sourcepointer-1}{}\fi
\noindent\textit{Source reference.} Original lemma 01XB: sources/stacks-chow/tags.txt, source lines 2512--2512. Cross-chapter standard prerequisite invoked once at proof lines 114-115 for the vanishing H\textasciicircum{}1(U, I) = 0 on an affine open; it supplies that vanishing only, while the local-ring comparison that concludes the proof is retained in the counted span.\par\medskip
\ifdefined\hypertarget\hypertarget{sourcepointer-2}{}\fi
\noindent\textit{Source reference.} Original lemma 01I1: sources/stacks-chow/tags.txt, source lines 1962--1962. Cross-chapter standard prerequisite invoked once at proof line 119 for the existence of the canonical morphism (U, A|\_U) -> Spec(A); the two commutative diagrams and the local-ring conclusion are retained in the counted span.\par\medskip
\begin{lemma}
\label{lemma-first-order-thickening}
Let $X$ be a scheme over a base $S$. Consider a short exact sequence
$$
0 \to \mathcal{I} \to \mathcal{A} \to \mathcal{O}_X \to 0
$$
of sheaves on $X$ where $\mathcal{A}$ is a sheaf of
$f^{-1}\mathcal{O}_S$-algebras,
$\mathcal{A} \to \mathcal{O}_X$ is a surjection
of sheaves of $f^{-1}\mathcal{O}_S$-algebras, and $\mathcal{I}$ is its kernel.
If
\begin{enumerate}
\item $\mathcal{I}$ is an ideal of square zero in $\mathcal{A}$, and
\item $\mathcal{I}$ is quasi-coherent as an $\mathcal{O}_X$-module
\end{enumerate}
then $X' = (X, \mathcal{A})$ is a scheme and $X \to X'$ is a first
order thickening over $S$. Moreover, any first order thickening over
$S$ is of this form.
\end{lemma}

\begin{proof}
It is clear that $X'$ is a locally ringed space. Let $U = \Spec(B)$
be an affine open of $X$. Set $A = \Gamma(U, \mathcal{A})$. Note that
since $H^1(U, \mathcal{I}) = 0$ (see Cohomology of Schemes, Lemma
\ref{coherent-lemma-quasi-coherent-affine-cohomology-zero})
the map $A \to B$ is surjective. By assumption the kernel
$I = \mathcal{I}(U)$ is an ideal of square zero in the ring $A$.
By
Schemes, Lemma \ref{schemes-lemma-morphism-into-affine}
there is a canonical morphism of locally ringed spaces
$$
(U, \mathcal{A}|_U) \longrightarrow \Spec(A)
$$
coming from the map $B \to \Gamma(U, \mathcal{A})$. Since this morphism
fits into the commutative diagram
$$
\xymatrix{
(U, \mathcal{O}_X|_U) \ar[d] \ar[r] & \Spec(B) \ar[d] \\
(U, \mathcal{A}|_U) \ar[r] & \Spec(A)
}
$$
we see that it is a homeomorphism on underlying topological spaces.
Thus to see that it is an isomorphism, it suffices to check it induces
an isomorphism on the local rings.
For $u \in U$ corresponding to the prime $\mathfrak p \subset A$
we obtain a commutative diagram of short exact sequences
$$
\xymatrix{
0 \ar[r] &
I_{\mathfrak p} \ar[r] \ar[d] &
A_{\mathfrak p} \ar[r] \ar[d] &
B_{\mathfrak p} \ar[r] \ar[d] & 0 \\
0 \ar[r] &
\mathcal{I}_u \ar[r] &
\mathcal{A}_u \ar[r] &
\mathcal{O}_{X, u} \ar[r] & 0.
}
$$
The left and right vertical arrows are isomorphisms because
$\mathcal{I}$ and $\mathcal{O}_X$ are quasi-coherent sheaves.
Hence also the middle map is an isomorphism. Hence every point
of $X' = (X, \mathcal{A})$ has an affine neighbourhood and $X'$ is a
scheme as desired.
\end{proof}

\end{document}
